\documentclass{article}

% Language setting
\usepackage[spanish]{babel}

\usepackage[letterpaper,top=2cm,bottom=5cm,left=3cm,right=3cm,marginparwidth=1.75cm]{geometry}

% Useful packages
\usepackage{amsmath}
\usepackage{graphicx}
\usepackage[colorlinks=true, allcolors=blue]{hyperref}
\usepackage{amssymb}
\usepackage{fancyhdr}
\usepackage{multicol}
\pagestyle{fancy}
\usepackage{url}
\usepackage{tabularx}
\usepackage{booktabs}
\usepackage{float}
\usepackage{pifont}

\fancyhf{} % Limpiar encabezado y pie de página por defecto
\setlength{\headheight}{70pt}

% Definición del encabezado
\fancyhead[C]{%
    \begin{minipage}[c]{2.0cm}
        \centering
        \includegraphics[width=2.0cm]{img/uta.png}
    \end{minipage}%
    \hfill
    \begin{minipage}[c]{8cm}
        \centering
        \small {\textbf{UNIVERSIDAD TÉCNICA DE AMBATO}} \\[0.1cm]
        \scriptsize \textit{FACULTAD DE INGENIERÍA EN SISTEMAS ELECTRÓNICA E INDUSTRIAL} \\[0.1cm]
        \scriptsize \textbf{CARRERA SOFTWARE - INTERACCIÓN HUMANO-COMPUTADOR} \\[0.1cm]
        \scriptsize \textbf{CICLO ACADÉMICO: AGOSTO  - DICIEMBRE 2026}
    \end{minipage}%
    \hfill
    \begin{minipage}[c]{2.0cm}
        \centering
        \includegraphics[width=2.0cm]{img/fisei.png}
    \end{minipage}%
}

\begin{document}
\begin{center}
    {\large \textbf{INFORME DE TAREA}} \\
\end{center}

\section{Portada}
\begin{flushleft}
\begin{tabular}{ l l }
    \textbf{Tema:} & Decisiones de diseño, Gestalt y guía de estilo -- Tutoría Fácil UTA \\
    \textbf{Unidad de Organización Curricular:} & Profesional. \\
    \textbf{Nivel y Paralelo:} & Quinto - B \\
    \textbf{Alumnos participante(s):} & Bejarano Masabanda Carlos Fabian\\
    & Guatemal Avilez Bryan Stalin \\
    & Molina Nuñez Karen Valeria\\
    & Solis Salinas Mario Tomas\\
    \textbf{Responsable del documento:} & Guatemal Avilez Bryan Stalin \\
    \textbf{Asignatura:} & Interacción Humano / Computador \\
    \textbf{Docente:} & Ing. Jose Ruben Caiza Caizabuano, Mg. \\
\end{tabular}
\end{flushleft}

\section{Informe de la Tarea}
\subsection{Objetivos}
\subsubsection{General}
Traducir las necesidades identificadas en el caso en decisiones de diseño
visibles y consistentes para Tutoría Fácil UTA. Para ello se aplican
metáforas, affordance, mapeo, manipulación directa, retroalimentación,
reducción de la carga cognitiva, leyes Gestalt y una guía mínima de estilo.

\subsubsection{Específicos}
\begin{itemize}
    \item Seleccionar metáforas y comportamientos de interacción coherentes
    con los modelos mentales de los estudiantes (E9).
    \item Aplicar y justificar al menos tres leyes Gestalt visibles en el prototipo.
    \item Definir tipografía, colores, componentes y reglas de consistencia
    para las cuatro pantallas.
\end{itemize}

\subsection{Instrucciones}
Las decisiones de diseño deben definirse antes de construir el prototipo.
Cada decisión debe resolver una necesidad del caso, vincularse con una
evidencia de E1 a E10 y aparecer después en el prototipo.

\subsection{Listado de equipos, materiales y recursos}
\textbf{Listado de equipos y materiales generales empleados en la tarea:}
\begin{itemize}
    \item Computador portátil y navegador web.
    \item Documento del examen y evidencias E1 a E10.
    \item FigJam, Figma Design y GitHub.
\end{itemize}

\textbf{TAC (Tecnologías para el Aprendizaje y Conocimiento) empleados en la tarea:}
\begin{itemize}
    \item[$\square$] Plataformas educativas
    \item[$\square$] Simuladores y laboratorios virtuales
    \item[$\boxtimes$] Aplicaciones educativas: FigJam y Figma Design
    \item[$\square$] Recursos audiovisuales
    \item[$\square$] Gamificación
    \item[$\boxtimes$] Inteligencia Artificial: apoyo en la redacción y el formato
    \item Otros (Especifique): \underline{GitHub}
\end{itemize}

\subsection{Desarrollo de la Tarea (Actividades a Desarrollar)}

\subsubsection{Metáforas}

Se emplean dos metáforas que los estudiantes ya reconocen (E9).

\textbf{Calendario.} Representa las fechas y la disponibilidad del docente.
Aprovecha un modelo mental conocido, porque las personas asocian
naturalmente un calendario con fechas y planificación.

\textbf{Tarjeta de tutoría.} Agrupa en un solo objeto el docente, la
asignatura, la fecha, la hora, la modalidad y el estado. Funciona como una
cita en una agenda física: el estudiante la consulta en \textit{Mis
tutorías} sin buscarla entre mensajes (E5).

\subsubsection{Affordance}

Los horarios disponibles tienen una apariencia claramente seleccionable:
borde de color, área táctil de al menos $44 \times 44$\,px y texto
explícito.

\begin{center}
\textbf{09:00 -- Disponible}
\end{center}

Los horarios ocupados aparecen en gris, con texto, y no se pueden accionar.
De esta forma se evitan los intentos de reservar horarios ya ocupados
observados en E3.

\begin{center}
\textbf{11:00 -- Ocupado}
\end{center}

\subsubsection{Mapeo}

La relación entre acción y resultado es directa. Cuando el estudiante toca

\begin{center}
\textbf{12:00 -- Disponible},
\end{center}

el mismo elemento cambia a

\begin{center}
\textbf{12:00 -- Seleccionado \ding{51}}
\end{center}

y el botón \textit{Continuar} se habilita justo debajo. El resultado se
muestra en el mismo elemento sobre el que se actuó, y la acción siguiente
aparece en la dirección natural de lectura.

\subsubsection{Manipulación directa}

El estudiante selecciona y modifica el horario actuando directamente sobre
el elemento visible, sin escribir ninguna hora. Para corregir, basta con
tocar otro horario. En la reprogramación, la acción \textit{Reprogramar}
abre el mismo calendario con la cita actual marcada, de modo que se usa el
mismo modelo de interacción. Esta decisión responde a E10.

\subsubsection{Retroalimentación}

La aplicación informa en todo momento qué está ocurriendo, siempre con texto:

\begin{itemize}
    \item \textbf{Carga:} \textit{Consultando horarios\ldots}, con el botón
    deshabilitado mientras se espera (E8).
    \item \textbf{Selección:} \textit{Horario seleccionado: jueves 15 de octubre, 12:00}.
    \item \textbf{Confirmación:} \textit{Tutoría CONFIRMADA}, junto con la
    tarjeta de la cita (E3, E5).
    \item \textbf{Error:} \textit{No pudimos cargar los horarios}, con el
    botón \textit{Intentar nuevamente} (E8).
\end{itemize}

\subsubsection{Reducción de la carga cognitiva}

El proceso se divide en etapas con un indicador de progreso
(\textit{Paso 1 de 2}): búsqueda, selección, revisión y confirmación.

\begin{itemize}
    \item Solo se muestran los horarios del día elegido y como máximo cinco
    días a la vez. Así se limitan las opciones simultáneas, en línea con la
    ley de Hick-Hyman y la capacidad de la memoria de trabajo (Miller).
    \item Se aplica el principio de reconocimiento sobre recuerdo: los datos
    de la cita se repiten antes y después de confirmar, así que el estudiante
    no tiene que memorizarlos.
\end{itemize}

Esta decisión se relaciona con E4 y E5.

\subsubsection{Riesgo cultural}

No se depende exclusivamente de colores, iconos o símbolos aislados. Por
ejemplo, en lugar de mostrar solo un símbolo de verificación se presenta:

\begin{center}
\textbf{\ding{51} Tutoría CONFIRMADA}
\end{center}

Cada horario incluye además su etiqueta textual: \textit{Disponible},
\textit{Ocupado} o \textit{Seleccionado}. Las fechas se escriben completas
(\textit{jueves 15 de octubre}) para evitar formatos ambiguos. Esto responde
a E2 y E9.

\subsubsection{Matriz de decisiones}

\begin{table}[H]
\centering
\caption{Decisiones de diseño aplicadas.}
\small
\begin{tabularx}{\textwidth}{p{2.5cm} X p{1.4cm} p{1.6cm}}
\toprule
\textbf{Concepto} & \textbf{Decisión aplicada} & \textbf{Evid.} & \textbf{Pantalla} \\
\midrule
Metáfora & Calendario para fechas y tarjeta para representar una tutoría. & E9 & P2, P4 \\
Affordance & Horarios disponibles seleccionables; ocupados en gris, con texto y deshabilitados. & E3, E9 & P2 \\
Mapeo & El horario seleccionado cambia de estado en el mismo lugar y \textit{Continuar} se habilita debajo. & E3, E9 & P2 \\
Manipulación directa & Selección y cambio actuando sobre el horario visible, también al reprogramar. & E10 & P2, P4 \\
Retroalimentación & Estados de carga, selección, confirmación y error con texto. & E3, E5, E8 & P2--P4 \\
Carga cognitiva & Pasos progresivos, pocas opciones a la vez y resumen visible. & E4, E5 & P2, P3 \\
Riesgo cultural & Ningún estado depende solo de iconos, símbolos o colores. & E2, E9 & Todas \\
\bottomrule
\end{tabularx}
\end{table}

\subsubsection{Aplicación de leyes Gestalt}

\paragraph{Proximidad.} Docente, fecha, hora y modalidad se agrupan dentro
de la misma tarjeta, con 8\,px de separación interna y 24\,px entre bloques
distintos. La cercanía permite entender que todos esos datos pertenecen a la
misma tutoría (E5).

\paragraph{Semejanza.} Los horarios que comparten estado tienen la misma
forma y estilo; los disponibles se distinguen de los ocupados por el estilo
y también por el texto (E3, E9). Además se mantiene la consistencia entre
botones principales, botones secundarios, tarjetas y mensajes de estado.

\paragraph{Figura--fondo.} Las tarjetas, los resúmenes y las confirmaciones
se destacan sobre el fondo gris mediante contraste, espacio y contenedores
blancos. Mientras el sistema carga, una capa oscurece el fondo. Así la
atención se dirige a los datos que el estudiante debe verificar.

\paragraph{Leyes complementarias.} Por \textit{continuidad}, el flujo es
vertical (docente $\rightarrow$ fecha $\rightarrow$ horario $\rightarrow$
Continuar). Por \textit{cierre}, el carrusel de días muestra el último día
cortado, lo que invita a desplazarse.

\subsubsection{Guía de estilo}

\begin{table}[H]
\centering
\caption{Guía mínima de estilo (frame móvil de $390 \times 844$\,px).}
\small
\begin{tabularx}{\textwidth}{p{3.2cm} X}
\toprule
\textbf{Elemento} & \textbf{Definición} \\
\midrule
Título principal & 24\,px, semibold. \\
Subtítulo & 18\,px, semibold. \\
Texto de cuerpo & 16\,px, regular. \\
Texto auxiliar & 14\,px, regular (nunca menor). \\
Acción principal & \texttt{\#0057B8} con texto blanco. \\
Fondo & \texttt{\#F7F8FA}; tarjetas en \texttt{\#FFFFFF}. \\
Texto principal & \texttt{\#1F2937}. \\
Éxito & \texttt{\#198754}, siempre acompañado de icono y texto (\textit{CONFIRMADA}). \\
Error & \texttt{\#DC3545}, siempre acompañado de un mensaje que explica la corrección. \\
Foco & Contorno de 3\,px, visible en todos los controles. \\
Componentes & Botón principal (48\,px de alto y ancho completo), botón secundario, campo con etiqueta, horario en tres variantes (disponible, ocupado y seleccionado), tarjeta de tutoría y mensaje de estado. \\
Etiquetas fijas & \textit{Ver horarios}, \textit{Mis tutorías}, \textit{Continuar}, \textit{Confirmar tutoría}, \textit{Cambiar horario}, \textit{Reprogramar} y \textit{Volver al inicio}. \\
Consistencia & La acción principal va siempre abajo; \textit{Volver} siempre arriba a la izquierda; se mantienen las mismas jerarquías y respuestas del sistema en las cuatro pantallas. \\
\bottomrule
\end{tabularx}
\end{table}

\subsubsection{Estándares de iure que fundamentan la guía}

La norma ISO 9241-11 orienta la calidad de uso en un contexto especificado
mediante la efectividad, la eficiencia y la satisfacción. La norma ISO
9241-210 orienta el proceso iterativo de diseño centrado en las personas.
Esta guía de estilo concreta ambas orientaciones en reglas consistentes y
verificables para el proyecto.

\subsection{Resultados obtenidos}
\begin{itemize}
    \item Se obtuvo una matriz de decisiones (Tabla 1) que relaciona los
    principios del parcial con necesidades concretas del caso, cada una con su
    evidencia y la pantalla donde se verifica.
    \item Se justificaron tres leyes Gestalt principales (proximidad,
    semejanza y figura-fondo) y dos complementarias (continuidad y cierre).
    \item Se definió una guía de estilo mínima (Tabla 2) que se aplica
    directamente en Figma y mantiene la consistencia de las cuatro pantallas.
\end{itemize}

\subsection{Habilidades blandas empleadas en la tarea}
\begin{itemize}
    \item[$\square$] Liderazgo
    \item[$\boxtimes$] Trabajo en equipo
    \item[$\boxtimes$] Comunicación asertiva
    \item[$\square$] La empatía
    \item[$\boxtimes$] Pensamiento crítico
    \item[$\boxtimes$] Flexibilidad
    \item[$\square$] La resolución de conflictos
    \item[$\boxtimes$] Adaptabilidad
    \item[$\boxtimes$] Responsabilidad
\end{itemize}

\subsection{Conclusiones}
\begin{itemize}
    \item Las metáforas de calendario y tarjeta aprovechan modelos mentales
    conocidos por los estudiantes y reducen la necesidad de explicar la interfaz.
    \item La retroalimentación textual y el mapeo directo disminuyen la
    ambigüedad al seleccionar y confirmar horarios, que era el principal
    problema del proceso actual (E3).
    \item Las leyes Gestalt organizan visualmente los elementos y reducen la
    carga perceptiva en una pantalla pequeña.
    \item La guía de estilo convierte criterios generales en reglas visuales
    concretas y reutilizables.
\end{itemize}

\subsection{Recomendaciones}
\begin{itemize}
    \item Aplicar los mismos componentes en las cuatro pantallas.
    \item No cambiar la etiqueta de una misma acción a lo largo del flujo.
    \item Verificar que todo estado con color incluya también texto.
    \item Comparar el prototipo final con la matriz de decisiones antes de la
    prueba cruzada.
\end{itemize}

\subsection{Referencias}
\begin{thebibliography}{9}

\bibitem{guia}
Universidad Técnica de Ambato, FISEI, ``Examen práctico integrador del primer
parcial: Diseño desde cero de un sistema web para reservar tutorías,''
Interacción Humano Computador, Ambato, Ecuador, 2026.

\bibitem{iso11}
International Organization for Standardization, \textit{ISO 9241-11:2018
Ergonomics of human-system interaction -- Part 11: Usability: Definitions and
concepts}. Ginebra, Suiza: ISO, 2018.

\bibitem{iso210}
International Organization for Standardization, \textit{ISO 9241-210:2019
Ergonomics of human-system interaction -- Part 210: Human-centred design for
interactive systems}. Ginebra, Suiza: ISO, 2019.

\end{thebibliography}

\subsection{Anexos}
Captura de la guía de estilo y de los componentes en Figma (insertar antes de la entrega).

\end{document}